\documentclass[11pt]{article}
\usepackage[T1]{fontenc}
\usepackage{lmodern,microtype}
\usepackage[margin=1in]{geometry}
\usepackage{amsmath,amssymb,amsthm,booktabs,hyperref}
\hypersetup{hidelinks,pdftitle={Projective-plane determinant bounds for sparse binary matrices},pdfauthor={John Tripodo}}
\emergencystretch=3em
\newtheorem{theorem}{Theorem}
\newtheorem{lemma}[theorem]{Lemma}
\newtheorem{corollary}[theorem]{Corollary}
\theoremstyle{remark}
\newtheorem{remark}[theorem]{Remark}
\newcommand{\tr}{\operatorname{tr}}
\title{Projective-plane determinant bounds\\for sparse binary matrices}
\author{John Tripodo\thanks{\copyright\ 2026 John Tripodo. Licensed under \href{https://creativecommons.org/licenses/by/4.0/}{CC BY 4.0}. ORCID: \href{https://orcid.org/0009-0006-9180-4953}{0009-0006-9180-4953}.}}
\date{21 October 2026}
\begin{document}
\maketitle

\begin{abstract}
For an integer $r\ge 3$ put $t=r-1$, $v=r^2-r+1$ and $D_r=r\,t^{(v-1)/2}$, the determinant magnitude of a projective-plane incidence matrix of order $t$ whenever such a plane exists. We prove that for every $3\le r\le 14$ and every
$n\times n$ matrix $A$ with entries in $\{0,1\}$ and $e$ ones,
\[
|\det A|\le \exp\!\bigl(n\,\ell_r+\alpha_r\,(e-rn)\bigr),\qquad \ell_r=\tfrac1v\log D_r,
\]
with an explicit constant $\alpha_r>0$. In particular $e\le rn$ implies $|\det A|\le D_r^{\,n/v}$, with no
restriction on individual row or column degrees. Equality holds exactly for direct sums of incidence matrices of
projective planes of order $r-1$ (up to row and column permutations). The bound is attained at positive orders divisible by $v$ for
$r\in\{3,4,5,6,8,9,10,12,14\}$, is never attained at positive order for
$r\in\{7,11\}$, and is attained at some positive order for $r=13$ if and only
if a projective plane of order $12$ exists. The case $r=3$, $|\det A|\le 24^{n/7}$, is the sparse
determinant conjecture of Bruhn and Rautenbach, recorded as Conjecture~1 by Araujo, Balogh and Wang.
The proof combines a degree-normalised Gram matrix, a quadratic Hermite majorant of $\log$, an edgewise
convex-tangent decomposition and cofactor peeling of low-degree lines. The only computer-assisted step is a
finite exact-integer certificate of an explicit two-variable inequality for each $r$, with rigorous coverage of every unbounded degree tail.
\end{abstract}

\section{Introduction}
A bound on the total number of ones permits arbitrarily dense individual rows
or columns. Thus total sparsity is distinct from a fixed row sum or simultaneous
local degree caps. The strengthened all-count inequality below is essential:
cofactor minors need not retain an original average-degree bound.

Ryser's fixed-order, fixed-total determinant bound has symmetric designs as
its equality matrices~\cite[Theorem~3]{Ry}. At projective-plane parameters
$n=r^2-r+1$ and $e=rn$, it gives the value $D_r$ below. Projective-plane blocks
and their determinant rates are classical; see also Scheinerman~\cite[Section~2]{S}.
The question here is a uniform inequality for every order and degree distribution.

Bruhn and Rautenbach~\cite{BR} obtained $6^{n/6}$ for at most $2n$ ones and
proposed the Fano-plane rate for at most $3n$ ones. Shitov~\cite{Sh} proved
$3^{(e-n)/4}$ for $e\ge n$. Araujo, Balogh and Wang~\cite[Theorem~3]{ABW}
improved this to $2^{(e-n)/3}$, giving $2^{2n/3}$ when $e\le3n$, and
recorded the Fano statement as Conjecture~1 in their Section~6.
Scheinerman~\cite[Theorem~7.1]{S} also treats unrestricted total-$kn$ sparsity,
with bound $c_k^n$, where
\[
 c_k=(k^2-1)^{(k-1)/(4k)}k^{1/(2k)}<\sqrt{k}.
\]
His results are not confined to fixed row sums. Semonsen~\cite{Se} gives
further fixed-row-sum estimates through linear programming.
Hu and Zuo~\cite[Theorem~4]{HZ} also give a more recent unrestricted
edge-count determinant estimate, alongside their separate work on $C_4$-free
perfect matchings. Their displayed determinant estimate does not give the
projective-plane constants proved here. Our result
supplies the projective-plane rate for the finite density range below and
identifies all equality matrices.

All logarithms are natural. Binary entries are regarded as real numbers when
taking norms and logarithms. We use the empty determinant convention
$\det(\varnothing)=1$. Blocks in a direct sum need not be isomorphic.

\begin{theorem}\label{thm:main}
Let $r$ be an integer with $3\le r\le 14$, let $n\ge0$ be an integer, put $t=r-1$, $R=r^2$, $v=R-t$, and define
\[
\gamma_r=\frac{v/t-\log(R/t)}{v^2},\qquad
\alpha_r=\frac1r-\frac1{2t}+\Bigl(r-\frac12\Bigr)\gamma_r,\qquad
\ell_r=\frac{\log r+\frac{v-1}{2}\log t}{v}.
\]
Then $\alpha_r>0$ (since $1/r-1/(2t)=(r-2)/(2rt)>0$ and $\gamma_r>0$, as $x-1>\log x$ at $x=R/t>1$), and
every $n\times n$ binary matrix $A$ with $e$ ones satisfies
\begin{equation}\label{eq:main}
|\det A|\le \exp\bigl(n\ell_r+\alpha_r(e-rn)\bigr).
\end{equation}
Equality holds if and only if, after independent permutations of rows and columns, $A$ is a direct sum of
incidence matrices of projective planes of order $r-1$ (the empty matrix included).
\end{theorem}

\begin{corollary}
For integers $3\le r\le 14$ and $e\le rn$, $|\det A|\le D_r^{\,n/v}$. The bound is attained, for every $n$ divisible by
$v$, when $r-1$ is a prime power, i.e.\ for $r\in\{3,4,5,6,8,9,10,12,14\}$. For $r=7$ and $r=11$ it is never
attained at positive order, since there are no projective planes of orders $6$ (Bruck--Ryser~\cite{BruckRyser}) and $10$ (Lam, Thiel and Swiercz~\cite{LTS}).
For $r=13$ it is attained at some positive order if and only if a projective plane of order $12$ exists.
Moreover, $e<rn$ or $v\nmid n$ forces strict inequality in the total-budget bound.
\end{corollary}

The existence classifications in this corollary are external design-theoretic
facts. In particular, nonattainment does not by itself prove a uniform
improvement in the asymptotic exponential rate. No such improvement is claimed
here for $r=7,11$, or $13$.

\paragraph{Proof outline.}
Degree normalization gives a spectral cap without bounding any original degree.
A logarithmic majorant converts the determinant into weighted Gram moments.
A convex tangent decomposition couples the row and column deficits along actual
incidences. A finite exact scalar certificate, including the entire degree tail,
then gives the bound whenever all degrees exceed a small threshold. Cofactor
expansion handles the remaining degrees. Vanishing of every lost term at
equality forces actual isolated projective-plane blocks.

\begin{table}[t]\centering\small
\begin{tabular}{rrrrlr}\toprule
$r$ & $v$ & $D_r^{1/v}$ & $K$ & Anchor & $10^7\min T_r$ exceeds\\\midrule
3&7&1.574610&1&linear&17482\\
4&13&1.847235&1&linear&37649\\
5&21&2.089198&1&linear&18842\\
6&31&2.308410&1&linear&10371\\
7&43&2.510041&1&linear&6167\\
8&57&2.697612&1&linear&3894\\
9&73&2.873633&1&linear&2579\\
10&91&3.039955&1&quadratic&1999\\
11&111&3.197992&2&linear&1262\\
12&133&3.348848&2&cubic&1035\\
13&157&3.493405&2&cubic&772\\
14&183&3.632379&2&cubic&588\\\bottomrule
\end{tabular}
\caption{Finite-window certificate, excluding $(r,r)$. The final column gives
strict rational bounds with denominator $10^7$. Rate decimals are illustrative
only. All tail-box and peeling-margin lower bounds exceed $3/100$.}
\label{tab:cert}
\end{table}

\section{The normalised Gram matrix and the vertex deficit}
Fix $r$ and drop the subscript. Let $A$ be nonsingular with all row degrees $d_i\ge1$ and column degrees $c_j\ge1$.
Put $M_{ij}=A_{ij}/\sqrt{d_ic_j}$ and $G=R\,MM^{\top}$.

\begin{lemma}[contraction]\label{lem:contract}
The Euclidean operator norm satisfies $\|M\|\le1$; hence the eigenvalues of $G$ lie in $(0,R]$.
\end{lemma}
\begin{proof}
By Cauchy--Schwarz, $(\sum_jM_{ij}x_j)^2\le(\sum_jA_{ij}/d_i)(\sum_jA_{ij}x_j^2/c_j)=\sum_jA_{ij}x_j^2/c_j$.
Summing over $i$ gives $\sum_j x_j^2$. Positivity follows from $\det A\ne0$.
\end{proof}

Since $\det G=R^n(\det A)^2/(\prod d_i\prod c_j)$,
\begin{equation}\label{eq:logdet}
\log|\det A|^2=\log\det G+\sum_i\log\frac{d_i}{r}+\sum_j\log\frac{c_j}{r}.
\end{equation}

\paragraph{Hermite majorant.} Let $q(z)=a+bz-\gamma z^2$ with $b=1/t+2\gamma t$, $a=\log t-bt+\gamma t^2$. Then
$q(t)=\log t$, $q'(t)=1/t$, $q(R)=\log R$. More explicitly, for $z>0$,
\[
 q(z)-\log z=(z-t)^2(R-z)
   \int_0^\infty\frac{ds}{(s+t)^2(s+R)(s+z)}.
\]
This identity follows by subtracting the Hermite interpolant of $1/(s+z)$
inside the integral representation of $\log z$. The integral converges and is
strictly positive. Hence $\log z\le q(z)$ on $(0,R]$, with equality exactly
at $z\in\{t,R\}$. Applying this to every eigenvalue gives
$\log\det G\le na+b\tr G-\gamma\tr G^2$.

\paragraph{Moments.} Let $u_i=\sum_jA_{ij}/(d_ic_j)$ and $w_j=\sum_iA_{ij}/(d_ic_j)$. Then $\tr G=R\sum_iu_i=R\sum_jw_j$,
$\sum_id_iu_i=\sum_jc_jw_j=n$, and expanding $\tr(MM^\top)^2$ over index quadruples gives
\[
\tr G^2=R^2\Bigl(\sum_iu_i^2+\sum_jw_j^2-\sum_{A_{ij}=1}\frac1{d_i^2c_j^2}+4W\Bigr),\qquad W\ge0,
\]
where $Q=\sum_{A_{ij}=1}1/(d_i^2c_j^2)$ and $W=\sum 1/(d_id_{i'}c_jc_{j'})$ over $i<i'$, $j<j'$ with $A_{ij}=A_{ij'}=A_{i'j}=A_{i'j'}=1$. The equal-row terms contribute $\sum_i u_i^2$, the equal-column terms
contribute $\sum_jw_j^2$, and their intersection is subtracted once. Each
rectangle has four orientations. This accounts for every term and the factor $4$.
Since
$\frac1{d^2c^2}\le\frac12(\frac1{d^3c}+\frac1{dc^3})$ (difference $\frac{(d-c)^2}{2d^3c^3}$), we obtain the following inequality. The upper bound for $Q$ has the correct
sign: it gives a lower bound for $\tr G^2$, whose coefficient is $-\gamma$.
The nonnegative rectangle loss is discarded.
\begin{equation}\label{eq:sep}
\log|\det A|^2\le na+\sum_iF(d_i,u_i)+\sum_jF(c_j,w_j),\qquad
F(d,u)=\Bigl(\frac{bR}2+\frac{\gamma R^2}{2d^2}\Bigr)u-\gamma R^2u^2+\log\frac dr .
\end{equation}

\paragraph{The vertex deficit.} With
$F_0=F(r,1/r)=br/2-\gamma(R-r/2)$, define
\[
q_r(d,u)=F_0+\alpha(d-r)+\beta(du-1)-F(d,u),\qquad \beta=\frac r{2t}-\gamma\Bigl(R+\frac r2\Bigr).
\]
The constants $\alpha,\beta$ are exactly those making $q_r$ vanish to first order at the plane point
$(d,u)=(r,1/r)$, and $a+2F_0=2\ell$. Indeed
\[
 a+2F_0=\log t+\frac1t-v\gamma
       =\log t+\frac{\log(R/t)}v=2\ell.
\]
These are algebraic identities; a symbolic replayer is supplied as a control.
Summing~\eqref{eq:sep} with $\sum_id_i=\sum_jc_j=e$ and $\sum d_iu_i=\sum c_jw_j=n$ yields
\begin{equation}\label{eq:deficit}
\log|\det A|\le n\ell+\alpha(e-rn)-\frac12\Bigl(\sum_iq_r(d_i,u_i)+\sum_jq_r(c_j,w_j)\Bigr).
\end{equation}

\section{Edgewise convex tangents}
For fixed $d$, $q_r(d,\cdot)$ is a convex quadratic ($\partial_u^2q_r=2\gamma R^2>0$). For an anchor function
$x\mapsto A(x)$ with $A(1/r)=1/r$, the supporting tangent at $A(1/d)$ gives
$q_r(d,u)\ge q_r(d,A)+\partial_uq_r(d,A)(u-A)$ for all real $u$. A row vertex of degree $d$ has
$u=\frac1d\sum_{j\sim i}\frac1{c_j}$, so the right side splits over its $d$ incident ones. Doing the same for
columns, the total deficit in~\eqref{eq:deficit} is at least $\sum_{A_{ij}=1}T_r(1/d_i,1/c_j)$ with
\[
T_r(x,y)=g(x,y)+g(y,x),\qquad g(x,y)=x\bigl[q_r(\tfrac1x,A(x))+\partial_uq_r(\tfrac1x,A(x))\,(y-A(x))\bigr].
\]
Multiplying through by $x$ removes all singularities: $g$ is a polynomial in $x,y$ plus $x\log(rx)$.
The anchors are
\[
A_{\rm lin}(x)=\tfrac{2}{5r}+\tfrac35x,\quad A_{\rm quad}(x)=-\tfrac1{100}+\tfrac65x-x^2\ (r=10),\quad
A_{\rm cub}(x)=x-\tfrac23(x-\tfrac1r)^2-\tfrac32(x-\tfrac1r)^3 .
\]
At the plane point, $q_r=\partial_uq_r=0$, hence $T_r(1/r,1/r)=0$ exactly.
For the finite certificate we use the expanded expression
\begin{align}\label{eq:Texpanded}
T_r(x,y)={}&2\alpha+(F_0-r\alpha)(x+y)+x\log(rx)+y\log(ry)-bRxy\nonumber\\
&+\gamma R^2\left[2xy\{A(x)+A(y)\}-xA(x)^2-yA(y)^2
                              -\frac{xy(x^2+y^2)}2\right].
\end{align}
Set $x\log(rx)=0$ at $x=0$; the formula then extends continuously to the
closed nonnegative quadrant. This explicit nonsingular form also specifies
exactly what the certificate evaluates.

\begin{lemma}[edge certificate]\label{lem:edge}
For each $3\le r\le 14$, with the anchor and integer $K$ of Table~1, $T_r(1/d,1/c)>0$ for all integers
$d,c\ge K+1$ with $(d,c)\neq(r,r)$.
\end{lemma}
\begin{proof}[Exact finite certificate]
Put $N=12r$ and $H=1/(N+1)$. By symmetry, all integer reciprocal degree pairs
are covered by: (i) the points $(1/d,1/c)$ with $K+1\le d\le c\le N$;
(ii) the boxes $\{1/d\}\times[0,H]$ for $K+1\le d\le N$; and (iii) the box
$[0,H]^2$. Exclude the known zero $(r,r)$ from (i).

The standard-library program \texttt{verify\_family\_exact.py} evaluates
\eqref{eq:Texpanded} by integer interval arithmetic with scale $10^{50}$.
Every logarithm is bounded using a rational odd-power series and an explicit
tail. Appendix~\ref{app:exact} gives the arithmetic and tail enclosure.
Every listed point and box has a strictly positive lower endpoint; no box
requires subdivision. Table~\ref{tab:cert} gives conservative rational
finite-window bounds. Every mixed or double-tail lower bound exceeds $3/100$.
The whole infinite degree range is therefore covered, rather than sampled.
The plane point is zero by exact tangency, not by numerical enclosure.
\end{proof}

\section{Cofactor peeling}
\begin{lemma}\label{lem:peel}
Suppose~\eqref{eq:main} holds for all binary matrices of order $n-1$, and $A$ (order $n$) has a row or column
of degree $k$ with $\log k+\alpha(r-k)<\ell$. Then~\eqref{eq:main} holds strictly for $A$.
\end{lemma}
\begin{proof}
Transpose if needed. Laplace expansion along the line gives $|\det A|\le\sum_{j}|\det A^{(j)}|$ over its $k$
nonzero positions; $A^{(j)}$ has order $n-1$ and at most $e-k$ ones. By induction and $\alpha>0$,
$|\det A|\le k\exp\bigl((n-1)\ell+\alpha(e-k-r(n-1))\bigr)$, which is strictly below the right side
of~\eqref{eq:main} by the hypothesis.
\end{proof}
The certified margins $\ell-\alpha(r-k)-\log k$ are positive for $k=1$ (all $r$) and for $k=2$ when
$11\le r\le14$; the smallest is $0.0352$ ($r=11$, $k=2$).

\section{Proof of Theorem~\ref{thm:main}}
Induct on $n$; the empty matrix gives equality and a singular matrix is trivial. A nonsingular matrix has
every line degree $\ge1$. If some line has degree $k$ with $1\le k\le K$, Lemma~\ref{lem:peel} gives strict
inequality. Otherwise all degrees are $\ge K+1$, and
\eqref{eq:deficit} with Lemma~\ref{lem:edge} gives~\eqref{eq:main}.

\section{Equality and actual incidence blocks}
Suppose equality holds at positive order. The matrix is nonsingular and
peeling is strict, so every degree exceeds $K$. Every incidence must have
$T_r=0$; hence all row and column degrees are $r$, and $G=AA^\top$.
Equality in the Hermite step puts the spectrum of $G$ in $\{t,R\}$.
Equality in the discarded rectangle term forces $W=0$, since $\gamma>0$ and
all rectangle weights are positive. Distinct rows therefore meet in at most
one column.

The symmetric matrix $C=AA^\top-tI$ is binary with diagonal entries one and
\[
 C^2=vC,\qquad C\mathbf1=v\mathbf1.
\]
The first identity follows from its spectrum $\{0,v\}$, and the second from
$r$-regularity on both sides. Each row of $C$ has $v$ ones. If $C_{ij}=1$,
the dot product of rows $i$ and $j$ is $v$, so those binary rows are identical.
If $C_{ij}=0$, their dot product is zero. Since $C_{ii}=1$, this partitions the
rows into classes of exactly $v$ elements, with $C$ a direct sum of $J_v$.

Rows from different classes have Gram entry zero; a column cannot meet two
classes. A class of $v$ rows has $rv$ incidences, while each associated column
has degree $r$. It therefore has exactly $v$ associated columns. Independent
row and column permutations put $A$ into square nonsingular blocks $B$ with
\[
 BB^\top=tI+J_v,\qquad B\mathbf1=B^\top\mathbf1=r\mathbf1.
\]
Since $BJ_v=J_vB=rJ_v$, we also have
$B^\top B=B^{-1}(tI+J_v)B=tI+J_v$. Thus each block is the incidence matrix
of a symmetric $2$-$(v,r,1)$ design, equivalently a projective plane of order
$t$. This derives actual isolated blocks, rather than only their spectral
parameters.

Conversely such a block has
$|\det B|^2=R\,t^{v-1}$, so $|\det B|=D_r$, and it has $rv$ ones.
Determinants multiply and the affine exponent adds under block sums.
The empty sum gives the order-zero equality. This completes the theorem.

\section{Scope and limitations}
\begin{remark}
With these anchors the edge inequality fails at $(3,3)$ for $r=15,16$, while degree-$3$ cofactor peeling is not
available there; the method as stated stops at $r=14$. For $r\le6$ a simpler separable argument (bounding each
$q_r(d,u)$ directly) already suffices; it fails for $r\ge7$ near $d\approx2r$.
\end{remark}
The certificate package contains a fresh exact-integer scalar replayer and a
separate replay of the earlier outward-rounded interval certificate.
Finite-field constructions for $q\in\{2,3,4,5,7,8,9,11,13\}$ verify
$AA^\top=qI+J$ and the determinant by exact Bareiss elimination. Symbolic
identities are checked for general $r$. Exhaustive small-matrix numerical
checks and hill-climbing searches are controls, not proof premises.
No bounded matrix census establishes an unbounded matrix-order quantifier.

Density $16$ has a separate proof using additional positive spectral terms
and local resource allocation. A separate technical companion covers the unrestricted inequality, actual-plane
equality classification, and literal-real-entry endpoints. That separate
result neither enlarges the density range of Theorem~\ref{thm:main} nor
establishes the density-$15$ member or an all-density theorem.
The classical projective-plane examples and the total-sparsity antecedents
above are not claimed as new.

\paragraph{AI assistance.}
OpenAI tools assisted mathematical exploration, proof engineering, computational checking, literature research, and manuscript preparation. The author is responsible for the mathematical statements, proofs, source comparisons, and final text.

\appendix
\section{Exact certificate arithmetic}\label{app:exact}
The only computational obligation in the proof is the explicit scalar lemma
and the finitely many peeling margins. The supplied verifier uses only the
Python standard library; all acceptance decisions compare integers.

For rational $q\in[1,2]$, set $z=(q-1)/(q+1)$. With $M=100$,
\[
 2\sum_{j=0}^{M-1}\frac{z^{2j+1}}{2j+1}
 \le\log q\le
 2\sum_{j=0}^{M-1}\frac{z^{2j+1}}{2j+1}
 +\frac{2z^{2M+1}}{(2M+1)(1-z^2)}.
\]
This is the positive-term $2\operatorname{arctanh}z$ series; bounding all
tail denominators below by $2M+1$ gives the geometric remainder. Inversion
and powers of $2$ reduce every positive rational argument to $[1,2]$; the
same formula encloses $\log2$.

Intervals are stored as integer endpoints divided by $S=10^{50}$.
Rational lower endpoints are rounded down and upper endpoints up.
Addition is exact on the grid. For multiplication, take the minimum and
maximum of the four endpoint products and round outward after division by
$S$. For division by an interval of one sign, enclose its reciprocal then
multiply. An interval containing zero is never inverted. Induction on the
expression therefore proves enclosure for every polynomial operation in
\eqref{eq:Texpanded}, including the coefficients $\gamma,\alpha,b,F_0$.

At a point $x=1/d$, enclose $x\log(rx)$ directly by the series. For a tail,
let $\widehat H\ge1/(N+1)$ be the grid-rounded upper endpoint. The verifier
checks $3r\widehat H<1$ and $\log3>1$ exactly. Thus the derivative
$\log(rx)+1$ is negative on $(0,\widehat H]$, so
\[
 x\log(rx)\in[\widehat H\log(r\widehat H),0]
       \quad(0\le x\le\widehat H).
\]
The lower endpoint is again enclosed by the rational series. The slightly
enlarged closed interval preserves every reciprocal tail point; neither
Euler's number nor a transcendental minimizer enters the program.

The program evaluates every unordered finite pair, $N-K$ mixed boxes, and
the single double-tail box. Every retained lower endpoint must be positive.
A timeout, an uncertain sign, or an unresolved box cannot count as a pass.
Three deliberately altered configurations must be rejected: linear anchors
at $r=10$ and $r=12$, and omitted degree-two peeling at $r=13$. Rejection
means not certified; it is not in itself a determinant counterexample.
The exact zero at the plane point remains an algebraic identity.

The complete receipt records scale, series length, source SHA-256, pair
counts, all lower bounds, and rejected controls. Standard and optimized
Python runs agree byte for byte. A separate arithmetic review checks interval
containment against exact fractions and derives \eqref{eq:Texpanded}
symbolically from the original tangent definition. The README supplies the
build commands and the frozen artifact manifest.

\section{Formal statement correspondence}
The companion Lean formalization provides the following statement correspondence.
In namespace
\texttt{SparseComponent.TotalFamily}:
\begin{itemize}
\item \texttt{determinant\_total\_family} proves the all-count integer bound
for $3\le r\le14$ with only binary entries as a matrix hypothesis.
\texttt{determinant\_total\_family\_budget} specializes the budget, and
\texttt{literal\_total\_family} states the explicit design base.
\item \texttt{determinant\_total\_family\_equality\_iff} and its budget
counterpart identify equality through actual row-overlap components and
their projective-plane incidence property.
\item \texttt{determinant\_real\_total\_family} is a separately proved
literal-real-entry interface. Equality is expressed on a canonical integer
binary representative proved entrywise identical to the real matrix, with
matching determinant and total count.
\item \texttt{determinant\_family\_order\_multiple} proves
$n=(r^2-r+1)\,\#\{\text{actual components}\}$ at equality. Its
nondivisibility counterpart gives strictness.
\end{itemize}
The formal scalar route uses exact rational kernel-reduced gates through
$3r$ and an analytic reciprocal tail. It is distinct from the $12r$ cover
in this paper. Concrete endpoints discharge every scalar premise; a generic
arbitrary-density interface retaining such premises is not an all-density
result.

The formalization uses only the standard axioms \texttt{propext}, \texttt{Classical.choice}, and \texttt{Quot.sound}. External projective-plane existence and nonexistence results are not formalized. Source versions, dependency pins, and reproducible build instructions are supplied with the companion code.

\medskip
\begin{thebibliography}{9}
\bibitem{ABW} I.~Araujo, J.~Balogh, and Y.~Wang,
Maximum determinant and permanent of sparse $0$--$1$ matrices,
\emph{Linear Algebra Appl.} \textbf{645} (2022), 194--228.
\href{https://doi.org/10.1016/j.laa.2022.03.020}{doi:10.1016/j.laa.2022.03.020}.
\bibitem{BR} H.~Bruhn and D.~Rautenbach,
Maximal determinants of combinatorial matrices,
\emph{Linear Algebra Appl.} \textbf{553} (2018), 37--57.
\href{https://arxiv.org/abs/1711.09935}{arXiv:1711.09935}.
\bibitem{BruckRyser} R.~H.~Bruck and H.~J.~Ryser,
The nonexistence of certain finite projective planes,
\emph{Canad. J. Math.} \textbf{1} (1949), 88--93.
\href{https://doi.org/10.4153/CJM-1949-009-2}{doi:10.4153/CJM-1949-009-2}.
\bibitem{HZ} Y.~Hu and Y.~Zuo,
An upper bound on the number of perfect matchings in $C_4$-free balanced bipartite graphs,
\emph{Linear Algebra Appl.} \textbf{752} (2027), 171--184; available online 2026.
\href{https://doi.org/10.1016/j.laa.2026.08.017}{doi:10.1016/j.laa.2026.08.017}.
\bibitem{LTS} C.~W.~H.~Lam, L.~Thiel, and S.~Swiercz,
The non-existence of finite projective planes of order 10,
\emph{Canad. J. Math.} \textbf{41} (1989), 1117--1123.
\href{https://doi.org/10.4153/CJM-1989-049-4}{doi:10.4153/CJM-1989-049-4}.
\bibitem{Ry} H.~J.~Ryser,
Maximal determinants in combinatorial investigations,
\emph{Canad. J. Math.} \textbf{8} (1956), 245--249.
\href{https://doi.org/10.4153/CJM-1956-028-4}{doi:10.4153/CJM-1956-028-4}.
\bibitem{S} D.~Scheinerman,
Maximal determinants of sparse zero-one matrices,
\href{https://arxiv.org/abs/1902.09644}{arXiv:1902.09644} (2019).
\bibitem{Se} J.~Semonsen,
An asymptotic determinant bound on $0$--$1$ matrices with fixed row sums,
\href{https://arxiv.org/abs/2002.02482}{\mbox{arXiv:2002.02482}} (2020).
\bibitem{Sh} Y.~Shitov,
On the determinant of a sparse $0$--$1$ matrix,
\emph{Linear Algebra Appl.} \textbf{554} (2018), 49--50.
\href{https://doi.org/10.1016/j.laa.2018.05.019}{doi:10.1016/j.laa.2018.05.019}.
\end{thebibliography}
\end{document}
